\chapter{Topology Optimization for Turbulent Flows: Results and Cross-Code CFD Verification}
\label{ch:benchmark_results}

\section{Introduction}
With the mathematical framework for turbulent topology optimization established in Chapter~\ref{ch:topology_optimization_methodology}, this chapter presents the resulting optimized topologies and rigorously evaluates their fluid dynamic performance. The objective is not merely to demonstrate that turbulent optimization yields different geometries than laminar assumptions, but to verify the physical credibility of the Spalart-Allmaras (SA) driven designs through independent computational fluid dynamics (CFD) analysis. 

To achieve this, a three-stage evaluation pipeline is utilized. First, the turbulent topology optimization is executed within the custom FEniCS framework, analyzing the objective function convergence, fluid routing, and eddy viscosity distributions. Second, the optimized continuous design fields are extracted as sharp, discrete geometries and exported to Ansys Fluent. These geometries are re-analyzed using the identical SA model to quantify cross-code consistency and assess the impact of geometry extraction. Finally, the designs are evaluated in Ansys using the higher-fidelity SST $k-\omega$ (and standard $k-\epsilon$) models to determine the sensitivity of the performance metrics to the underlying turbulence closure, particularly in regimes characterized by high Dean numbers and flow separation.

\section{Boundary Conditions and Evaluation Metrics}
\label{sec:boundary_conditions_metrics}
To accurately capture the physics of internal turbulent manifolds, the boundary conditions for all benchmark cases strictly utilize uniform inlet velocity profiles combined with a prescribed turbulence intensity (TI), diverging from the fully developed parabolic profiles typical of laminar topology optimization. Furthermore, to prevent the optimizer from exploiting artificial shear layers generated directly at the boundary constraints, non-design fluid regions ($\gamma = 1$) are enforced at all inlets and outlets. This allows the turbulent flow to naturally develop before interacting with the actively optimized design domain.

The primary performance metric is the total power dissipation (or equivalently, the total pressure drop for incompressible, fixed-flow-rate systems). Direct comparison of the raw objective function $J$ between FEniCS and Ansys requires careful interpretation. The FEniCS objective integrates both physical viscous dissipation and the artificial Brinkman penalty across a diffuse field. Conversely, Ansys evaluates purely physical dissipation on a body-fitted mesh bounded by sharp walls. Therefore, while the FEniCS objective history demonstrates the internal convergence of the optimizer, the ultimate validation metric relies on the post-processed hydraulic power loss calculated on the extracted Ansys geometries.

Geometry extraction inherently introduces uncertainties. The continuous design variable $\bar{\gamma}$ is thresholded at a specific iso-value to define the solid-fluid interface. Any discrepancies in dissipation between FEniCS and Ansys stem from a combination of discretization differences, solver architecture, and this thresholding process itself.

\section{Benchmark 1: The Diffuser}
\label{sec:benchmark_diffuser}
The diffuser serves as the foundational benchmark to evaluate separation control in a purely expanding flow, isolating shear layer dynamics from curvature-induced secondary flows. 

\subsection{FEniCS-SA Optimization Results}
The optimization relies on the frozen turbulence adjoint to suppress recirculation zones and minimize the pressure drop across the expansion. [Insert text analyzing the objective function plot: note how rapidly the objective drops initially before asymptoting]. 

Figure [X] illustrates the final projected topology, alongside the velocity magnitude and eddy viscosity fields. [Insert critical analysis of the SA field: Does the optimizer successfully maintain attached flow? Does the eddy viscosity field indicate an over-prediction of diffusion that artificially stabilizes the flow?]

\subsection{Cross-Code Verification and Model Sensitivity}
The thresholded geometry was imported into Ansys Fluent. Table [X] details the power dissipation evaluated across the different solvers and turbulence models. 

Comparing FEniCS-SA to Ansys-SA reveals a relative difference of [X]\%. This discrepancy highlights the impact of moving from a Brinkman-penalized diffuse interface to a wall-resolved, body-fitted boundary. When analyzed with the SST $k-\omega$ model, the total dissipation shifts by [X]\%. [Insert critical analysis: Does the SST model predict flow separation that the SA model missed during optimization? If the SA design fails to perform well under SST conditions, acknowledge this openly as a limitation of the one-equation closure in expanding adverse pressure gradients].

\section{Benchmark 2: The Double Pipe}
\label{sec:benchmark_double_pipe}
The double pipe configuration introduces complex merging and splitting dynamics, assessing the optimizer's ability to mitigate internal junction losses and mixing dissipation without the complications of global curvature.

\subsection{FEniCS-SA Optimization Results}
[Insert analysis of FEniCS results. Focus on the central junction. How does the optimizer route the flow to avoid stagnation zones? Reference the $\nu_t$ field to show where turbulence production is localized.]

\subsection{Cross-Code Verification and Model Sensitivity}
Table [X] presents the comparative CFD metrics for the double pipe geometry. [Insert analysis of the Ansys results. Evaluate if the sharp geometry extraction preserved the smooth junction routing. Compare SA vs SST $k-\omega$: Are the mixing losses at the intersection heavily dependent on the chosen turbulence model?]

\section{Benchmark 3: The Pipe Bend}
\label{sec:benchmark_pipe_bend}
To strictly evaluate the framework's performance in curvature-dominated regimes, the pipe bend benchmark proposed by Alexandersen et al. (2026) is utilized. This case incorporates explicit non-design regions and forces the fluid through a severe geometric deflection, triggering strong streamline curvature.

\subsection{Dean Number and Theoretical Limits}
As established in Chapter~\ref{ch:turbulence_modeling}, the standard Boussinesq approximation utilized by the SA model struggles to accurately capture the anisotropic stress tensors generated by strong secondary flows. The operational Dean number ($De$) for this benchmark is calculated as [Insert Value], based on the characteristic inlet diameter and the mean radius of curvature generated by the optimizer.

\subsection{Optimization and Re-analysis}
[Insert analysis of the FEniCS TO results]. 

Re-analysis in Ansys provides critical insight into the physical limits of the SA-optimized design. As shown in Table [X], the performance gap between Ansys-SA and Ansys SST $k-\omega$ is [Insert Result]. [If the gap is large: This discrepancy confirms the theoretical limitations outlined in Chapter~\ref{ch:turbulence_modeling}. The SA-driven optimizer, blind to curvature-induced anisotropy, generates a topology that underperforms when subjected to higher-fidelity turbulence models. This is not a failure of the numerical implementation, but a demonstrated physical limit of utilizing a one-equation linear eddy viscosity model for highly curved manifolds.]

\section{Benchmark 4: The U-Bend (Advanced Separation)}
\label{sec:benchmark_u_bend}
% Note: Include this section only if you successfully complete the Alexandersen U-bend case.
The U-bend benchmark forces a 180-degree flow reversal, combining intense curvature with inevitable flow separation at the inner radius. [Present the FEniCS results and Ansys comparisons. Focus on whether the optimizer managed to delay separation, and whether SST $k-\omega$ validates or rejects the SA-generated inner wall contour].

\section{Evaluation of the Coupled ``Semi-Frozen'' Adjoint}
\label{sec:semifrozen_evaluation}
While the frozen adjoint serves as the primary optimization engine due to its numerical tractability, an advanced ``semi-frozen'' adjoint was formulated to explicitly capture turbulence production sensitivities. 

However, evaluating this coupled formulation reveals significant practical limitations. As observed in Figure [X], while the coupled adjoint mathematically detects shear-driven turbulence, exact continuous adjoints for turbulence models are notoriously susceptible to numerical noise at domain boundaries. Even with non-design regions implemented, the coupled sensitivities frequently become dominated by artificial shear concentrations at fluid-solid transitions. Consequently, the optimizer tends to prioritize mitigating these localized mathematical artifacts rather than optimizing the global fluid routing. Therefore, despite its theoretical superiority, the frozen adjoint remains the more robust and reliable approach for the manifolds studied in this thesis.